% !TeX root = ../sections/02_exercises.tex
\subsection{2-C.}
\begin{exercise}
	\textup{(Heine--Borel の定理)}
	\(a,b\in\mathbb{R}\) が
	\[
	a<b
	\]
	をみたすとする。
	開区間の族
	\[
	\{U_\lambda\}_{\lambda\in\Lambda}
	\]
	が
	\[
	[a,b]\subset
	\bigcup_{\lambda\in\Lambda}U_\lambda
	\]
	をみたすとき，有限個の
	\[
	\lambda_1,\ldots,\lambda_n\in\Lambda
	\]
	が存在して
	\[
	[a,b]\subset
	\bigcup_{k=1}^{n}U_{\lambda_k}
	\]
	が成立することを証明せよ。
\end{exercise}

\begin{background}
	この問題は，閉区間 \([a,b]\) のコンパクト性を示す問題である。
	主張は，閉区間を開区間の族で覆うことができるならば，
	実はその中から有限個だけを選んでも閉区間全体を覆える，
	という内容である。
	
	このような有限個の部分族を
	有限部分被覆という。
	つまり，仮定は
	\[
	[a,b]\subset
	\bigcup_{\lambda\in\Lambda}U_\lambda
	\]
	であり，示すべきことは
	\[
	[a,b]\subset
	U_{\lambda_1}\cup\cdots\cup U_{\lambda_n}
	\]
	となる有限個の添字
	\(\lambda_1,\ldots,\lambda_n\) が取れることである。
	
	証明には区間縮小法を用いる。
	背理法で，\([a,b]\) が有限個の \(U_\lambda\) では覆えないと仮定する。
	そのうえで，区間を半分に分けていく。
	もし左右の半分がどちらも有限個で覆えるなら，
	元の区間も有限個で覆える。
	したがって，少なくとも一方は有限個では覆えない。
	このような半分を選び続けることで，
	有限個では覆えない閉区間の列を作る。
	
	区間の長さは半分ずつ小さくなるので，
	区間縮小法により，それらの共通部分にはただ一つの点が存在する。
	その点はもとの開被覆のどれか一つに含まれるため，
	十分小さい閉区間はその開区間一つに含まれる。
	これは「有限個では覆えない」とした選び方に反する。
\end{background}

\begin{strategy}
	背理法で証明する。
	すなわち，
	\[
	[a,b]
	\]
	は
	\[
	\{U_\lambda\}_{\lambda\in\Lambda}
	\]
	の有限個では覆えないと仮定する。
	
	まず
	\[
	I_1=[a,b]
	\]
	とおく。
	\(I_1\) を二等分して，
	\[
	I_1'
	=
	\left[
	a,\frac{a+b}{2}
	\right],
	\qquad
	I_1''
	=
	\left[
	\frac{a+b}{2},b
	\right]
	\]
	を考える。
	もし \(I_1'\) と \(I_1''\) がどちらも有限個の
	\(U_\lambda\) で覆えるならば，
	その有限個を合わせることで \(I_1=[a,b]\) も有限個で覆える。
	これは仮定に反する。
	
	したがって，少なくとも一方は有限個では覆えない。
	そのような方を選んで
	\[
	I_2
	\]
	とおく。
	
	同様に，すでに \(I_n\) が選ばれていて，
	\(I_n\) が有限個では覆えないとする。
	\(I_n\) を二等分し，そのうち有限個では覆えない方を
	\[
	I_{n+1}
	\]
	として選ぶ。
	
	この操作により，閉区間の列
	\[
	I_1\supset I_2\supset I_3\supset\cdots
	\]
	が得られる。
	さらに，各 \(I_n\) は有限個の \(U_\lambda\) では覆えず，
	区間の長さは
	\[
	|I_n|=\frac{b-a}{2^{n-1}}
	\]
	である。
	したがって
	\[
	|I_n|\to 0
	\]
	である。
	
	区間縮小法より，
	ある \(\alpha\in\mathbb{R}\) が存在して
	\[
	\bigcap_{n=1}^{\infty}I_n
	=
	\{\alpha\}
	\]
	となる。
	特に
	\[
	\alpha\in [a,b]
	\]
	である。
	
	仮定より
	\[
	[a,b]\subset
	\bigcup_{\lambda\in\Lambda}U_\lambda
	\]
	だから，ある \(\lambda_0\in\Lambda\) が存在して
	\[
	\alpha\in U_{\lambda_0}
	\]
	となる。
	\(U_{\lambda_0}\) は開区間なので，ある \(\varepsilon>0\) が存在して
	\[
	(\alpha-\varepsilon,\alpha+\varepsilon)
	\subset
	U_{\lambda_0}
	\]
	となる。
	
	一方，\(I_n\) の長さは \(0\) に収束し，さらに
	\(\alpha\in I_n\) である。
	したがって，十分大きい \(N\) に対して
	\[
	I_N
	\subset
	(\alpha-\varepsilon,\alpha+\varepsilon)
	\subset
	U_{\lambda_0}
	\]
	が成り立つ。
	
	これは \(I_N\) がただ一つの開区間
	\[
	U_{\lambda_0}
	\]
	で覆えることを意味する。
	したがって \(I_N\) は有限個の \(U_\lambda\) で覆える。
	これは，\(I_N\) を有限個では覆えないように選んだことに反する。
	
	よって背理法により，有限個の
	\[
	\lambda_1,\ldots,\lambda_n\in\Lambda
	\]
	が存在して
	\[
	[a,b]\subset
	\bigcup_{k=1}^{n}U_{\lambda_k}
	\]
	が成り立つ。
\end{strategy}